\section{评测结果}
\label{sec:results}

我们按问题而不是按研究来报告结果：控制器应当从模型手里拿走哪些决定（\S\ref{sec:rq-stop}），控制器采用验证器的证据标准之后模型还剩下什么（\S\ref{sec:fair}），模型读取原始日志时的表现（\S\ref{sec:rawlog}），由此形成的安全边界在哪里（\S\ref{sec:adversarial}），控制器的知识从哪里来、知识出错时如何失败（\S\ref{sec:mismatch}），以及效应能否在第二任务族上重复（\S\ref{sec:external}）。附录~\ref{app:tables} 给出每项研究的每一格。

\subsection{探测与停止是两个独立的决定}
\label{sec:rq-stop}
\label{sec:rq-rank}
由控制器选择探针同时改变了两件事：拿走了模型自己的选择权，并用一个排序取代它。分解研究用盲配置把两者分开（无论用哪个选择器，规划器收到的提示都相同）；停止研究再把“谁选探针”与“谁可以停止”交叉。\cref{fig:probe-stop} 展示后者。

\begin{figure*}[t]
  \centering
  \includegraphics[width=0.92\textwidth]{fig-probe-stop.pdf}
  \caption{停止研究：按谁选探针、谁可以停止给出的已验证完成率（每点 72 个回合）。虚线是事后运行的不含 LLM 的控制器（EIG/成本选择加控制器停止，种子和任务相同）。}
  \label{fig:probe-stop}
\end{figure*}

第一，排序本身有用。在规划器信息完全相同时，EIG/成本选择在 Qwen2.5-7B 上比随机选择高 $+0.306$ $[+0.14,+0.47]$（预先登记的终点），在其他每个会下结论的模型上高 $+0.26$ 到 $+0.35$。相反，拿走选择权对弱模型有利、对强模型不利：用随机选择代替模型自己的选择，Qwen2.5-7B 的完成率提高 $+0.556$，Qwen2.5-32B 和 72B 分别降低 0.222 和 0.278。向规划器展示排序对完成率的影响至多 $+0.06$。

第二，由控制器选探针并不能让模型下结论。在我们的提示和服务配置下，Llama-3.1-8B 在分解研究的 360 个回合中一个也没有完成：它的 3{,}389 次决策全部是请求下一个探针，输出均为合法 JSON。加上控制器停止后，它从 0 升到 0.917（$+0.917$ $[+0.79,+1.00]$，预先登记），此后控制器选择再增加 $+0.056$ $[-0.07,+0.18]$，没有可检测的完成率提升，但探针成本减半（2.35 对 4.78 个单位）。Qwen2.5-7B 两部分都需要：仅停止就把它从 0.083 提升到 0.694，选择再增加 $+0.222$ $[+0.08,+0.39]$。控制器事先无法知道某个模型缺的是哪一部分，所以两部分都要提供。

第三，后验停止会靠排除法下结论。在控制器选择加停止下，五个模型中有四个在同样的六个回合上失败：控制器通过排除其他可能得出正确的 DNS 命令与控制结论，但验证器因缺少确认特征而拒绝。在不启用控制器停止的配置中，自行决定何时停止的强模型会一直探测到确认特征出现。采用后验停止的不含 LLM 控制器以 2.35 个单位完成 0.917 的案例，而当控制器既选探针又停止时，四个模型的结果与它完全相同。

\subsection{模型还剩下什么}
\label{sec:fair}
上面的比较混在一起的是两件事：LLM 是否有帮助，以及用什么证据标准结束回合。v1.1 研究消除了这一混杂：它给不含 LLM 的控制器配上 \S\ref{sec:finish} 的确认条件停止，加入固定流程和目录顺序作为更简单的参照，并在两个任务族上以相同的种子、任务和预算运行所有 LLM 配置，全部开启结束守卫。其主要终点比较“Qwen2.5-72B 决定何时停止”与“不含 LLM 的确认条件停止”（97.5\% 区间）。\cref{tab:v11fair} 同时报告满足证据政策的回合数和原因正确的回合数。

\begin{table}[t]
  \caption{模型还剩下什么（v1.1）：满足证据政策的回合数 / 原因正确的回合数，分别在由 LLM 决定何时停止（“LLM”）、后验停止（“post.”）和确认条件停止（“conf.”）下。所有 LLM 行都使用控制器选择和结束守卫；种子、任务和预算与不含 LLM 的一行相同。}
  \label{tab:v11fair}
  \centering\small
  \begin{adjustbox}{max width=\columnwidth}% generated by paper/make_tables.py -- do not edit
\begin{tabular}{llrrrr}
\toprule
 & & \multicolumn{2}{c}{sec-triage (72)} & \multicolumn{2}{c}{sigma-triage (152)} \\
\cmidrule(lr){3-4}\cmidrule(lr){5-6}
Planner & Controller & verified & cost & verified & cost \\
\midrule
LLM-free & posterior stop & 66 & 2.35 & 152 & 2.20 \\
LLM-free & confirmation stop & 72 & 2.54 & 152 & 2.20 \\
LLM-free & fixed playbook, conf. & 66 & 2.90 & 152 & 2.20 \\
LLM-free & catalogue order, conf. & 65 & 4.93 & 151 & 2.96 \\
\addlinespace
Q-7B & guard, LLM stops & 66 & 7.17 & 126 & 3.22 \\
Q-7B & guard + posterior stop & 66 & 2.35 & 141 & 2.12 \\
Q-7B & guard + confirmation stop & 72 & 2.54 & 141 & 2.12 \\
\addlinespace
L-8B & guard, LLM stops & 0 & 10.00 & 1 & 9.78 \\
L-8B & guard + posterior stop & 66 & 2.35 & 152 & 2.20 \\
L-8B & guard + confirmation stop & 72 & 2.54 & 152 & 2.20 \\
\bottomrule
\end{tabular}
\end{adjustbox}
\end{table}

第一，采用验证器的标准后，控制器单独就能完成每一个案例：sec-triage 72/72、sigma-triage 152/152，成本分别为 2.54 和 2.20 个单位，在 sec-triage 上只比后验停止多 0.19 个单位。在 sec-triage 上自适应选择仍然重要（固定流程 66、目录顺序 65），在 sigma-triage 上则不然，任何顺序都至少完成 152 个中的 151 个。

第二，没有任何模型能在这个控制器之上增加完成率。Qwen2.5-72B 自行决定何时停止时，在 sec-triage 上 72/72 满足证据政策（$+0.000$ $[0.000, 0.000]$），在 sigma-triage 上 148/152（$-0.028$ $[-0.111, 0.000]$；这是按规则求差后在规则间的平均，逐回合为 $-4/152 = -0.026$），成本为 3.86 和 3.25 个单位，而控制器为 2.54 和 2.20。按照运行前写定的判读，我们没有测到 LLM 在停止判断上的增益。Qwen2.5-32B 和 Llama-3.1-70B 表现相同，三个大模型在 sigma-triage 上的所有未验证回合都是原因正确但缺少确认特征。由模型决定何时停止时，Llama-3.1-8B 只完成 0 和 1 个回合，Qwen2.5-7B 为 66 和 126 个。Qwen2.5-7B 在控制器停止下也有不足（152 个中 141 个），原因只有一个：在三条规则上，它在后验 0.794 到 0.796（略低于守卫阈值）时提出结束，守卫拒绝后它反复提出同样的结束，直到决策预算耗尽（11 个回合，每个被拒 9 到 11 次）。控制器停止在 0.8 以下不会触发，因此规划器可以让带守卫的控制器陷入僵局。v1.2 研究加入恢复规则：连续两次结束被拒后，控制器运行自己排名第一的探针。一个坚持在分数 0.5 时就结束的脚本规划器，没有恢复规则时让 sigma-triage 的全部 152 个回合卡死，有恢复规则时全部 152 个完成。

第三，证据政策是自动化与安全之间的取舍，而不是评分细节。我们的验证器只在某条被引用的观测带有只有该原因才会产生的特征时才接受结论，因此通过组合或排除得到的正确原因算作失败；\cref{tab:v11fair} 显示，若以“原因正确”为标准，后验停止同样能完成 72/72。v1.2 研究用不含 LLM 的控制器测量了其他政策（\cref{tab:v12policy}）。它新增 \emph{comp-triage}：第三个任务族，三个场景中攻击与三种合法活动使用相同工具，结果为概率分布，十二个原因中有五个没有单条独有观测；新增\emph{联合}验证器：只要全部观测在真实分布下对该原因的支持至少是其他每个具名原因的十倍即接受；并新增要求对具名原因、或同时对 \textit{other} 具有联合证据的停止规则。在 comp-triage 上，联合证据在 72 个回合中 59 个给出正确原因，单条特征规则只有 20 个，其余 52 个升级。代价以攻击者得益的方式付出：联合证据把 18 个攻击中的 4 个结为与之配对的合法活动；当真实原因不在候选集合中时，它在 sec-triage 的 48 个回合中有 24 个给出另一个需处置原因。要求证据比 \textit{other} 高十倍并没有改变这一点，三十倍也只降到 18。只有单条特征规则在所有设置下都没有给出错误原因，代价是大部分组合证据案例被升级；把它的比值降到 3，只是在同一前沿上移动（37 个正确原因、1 个漏判攻击）。部署时必须在这条前沿上选一个点；我们的选择偏向升级。

\begin{table*}[t]
  \caption{证据政策（v1.2，不含 LLM 的控制器，EIG/成本选择；计数）。“single” / “joint”：被单条特征验证器和联合识别验证器接受的回合数。sec, cause absent：真实原因从候选集合中移除。sec, confusion：每个原因的表行向另一个原因混合 0.25。comp-triage：三个场景，攻击与使用相同工具的合法活动配对，结果为概率分布，十二个原因中五个没有单条独有观测。}
  \label{tab:v12policy}
  \centering\small
  \begin{adjustbox}{max width=\textwidth}% generated by paper/make_tables.py -- do not edit
\begin{tabular}{lrrrrrrrrr}
\toprule
 & \multicolumn{2}{c}{sec, matched (48)} & \multicolumn{2}{c}{sec, cause absent (48)} & sec, confusion & \multicolumn{4}{c}{comp-triage (72; 18 attacks)} \\
\cmidrule(lr){2-3}\cmidrule(lr){4-5}\cmidrule(lr){6-6}\cmidrule(lr){7-10}
Stop rule & single & joint & wrong & escal. & verified & right cause & wrong & missed attack & escal. \\
\midrule
posterior (0.8) & 42 & 48 & 24 & 24 & 47 & 64 & 7 & 5 & 1 \\
single signature, $r{=}10$ & 48 & 48 & 0 & 48 & 0 & 20 & 0 & 0 & 52 \\
single signature, $r{=}3$ & 48 & 48 & 0 & 48 & 0 & 37 & 2 & 1 & 33 \\
joint, named causes, $r{=}10$ & 42 & 48 & 24 & 24 & 47 & 59 & 6 & 4 & 7 \\
joint, incl.\ \textit{other}, $r{=}10$ & 42 & 48 & 24 & 24 & 47 & 59 & 6 & 4 & 7 \\
joint, incl.\ \textit{other}, $r{=}30$ & 42 & 48 & 18 & 30 & 18 & 55 & 2 & 1 & 15 \\
\bottomrule
\end{tabular}
\end{adjustbox}
\end{table*}

第四，停止失败在一定程度上是提示造成的。回放 60 个证据压倒性或已无合法探针的归档请求，Llama-3.1-8B 每次都返回合法、未截断的 JSON；在原始提示、带结束示例的提示和明确说明何时结束的提示下，分别有 60、60 和 58 次请求探针。在 \hesp{} 内的在线回合中，明确规则让它在 48 个中完成 8 个，另外两个变体为 0；对 Qwen2.5-7B，两个变体都把 \hesp{} 内的完成数从 43 提高到 48/48。为检验基线是否只是缺少清楚的提示，v1.2 研究在开发种子上按事先写定的规则选出对基线帮助最大的变体（说明何时结束、给出示例、禁止重复探针的提示），再用两种提示分别运行所有配置（\cref{tab:v12fair}）。更清楚的提示让 Qwen2.5-7B 的 Memory-only 完成数翻倍（48 个中从 8 到 16），对 ReAct 式和两个 Llama 基线没有帮助。控制器的优势依然存在：在更清楚的提示下，确认条件停止比 Memory-only 在 Qwen2.5-7B 上高 $+0.667$ $[+0.46,+0.85]$，在 Llama-3.1-8B 上高 $+1.000$ $[+1.000,+1.000]$（预先登记，97.5\% 区间）。在带守卫的 \hesp{} 内，更清楚的提示让 Llama-3.1-8B 在 48 个回合中自行完成 32 个，原始提示下为 0。

\begin{table}[t]
  \caption{同样清楚的提示（v1.2，sec-triage，测试种子；已验证回合数，共 48）。“clearer”：在开发种子上为基线选出的变体。}
  \label{tab:v12fair}
  \centering\small
  \begin{adjustbox}{max width=\columnwidth}% generated by paper/make_tables.py -- do not edit
\begin{tabular}{lrrrr}
\toprule
 & \multicolumn{2}{c}{Qwen2.5-7B} & \multicolumn{2}{c}{Llama-3.1-8B} \\
\cmidrule(lr){2-3}\cmidrule(lr){4-5}
Configuration & original & clearer & original & clearer \\
\midrule
ReAct-style & 2 & 0 & 0 & 0 \\
Memory-only & 8 & 16 & 0 & 0 \\
\hesp{}, guard, LLM stops & 44 & 48 & 0 & 32 \\
\hesp{}, guard + confirmation stop & 48 & 48 & 48 & 48 \\
\bottomrule
\end{tabular}
\end{adjustbox}
\end{table}

\subsection{读取原始日志}
\label{sec:rawlog}
到目前为止，所有结果都向控制器提供离散观测。v1.2 研究在 sec-triage 上去掉了这一简化：每个探针返回日志文本，由解析器把它转换成账本使用的结果。我们比较按文档日志格式编写的规则解析器与被服务的模型作为解析器，在三种日志条件下的表现：文档格式；\emph{改写}格式，即用改写过的句子陈述同样的事实并加入无关日志行，相当于日志管道变更之后；\emph{注入}格式，即加入一个攻击者控制的字段，用探针自身的词表声称一个良性读数。规划由不含 LLM 的控制器完成（EIG/成本、确认条件停止），模型唯一的作用是读取；验证器检查真实结果。\cref{tab:v12raw} 报告结果。

\begin{table}[t]
  \caption{读取原始日志（v1.2，sec-triage base 与 noise，每格 48 个回合；不含 LLM 的控制器加确认条件停止）。解析准确率按全部被解析的观测计算；注入被采纳：真实结果不同、而解析器返回了攻击者所声称标签的观测。}
  \label{tab:v12raw}
  \centering\small
  \begin{adjustbox}{max width=\columnwidth}% generated by paper/make_tables.py -- do not edit
\begin{tabular}{llrrrrrr}
\toprule
Parser & Log format & \multicolumn{2}{c}{parse accuracy} & \multicolumn{2}{c}{injection adopted} & verified (of 48) & escalated (of 48) \\
\cmidrule(lr){3-4}\cmidrule(lr){5-6}\cmidrule(lr){7-7}\cmidrule(lr){8-8}
 & & Q-7B & L-8B & Q-7B & L-8B & Q-7B / L-8B & Q-7B / L-8B \\
\midrule
structured (reference) & documented & -- & -- & -- & -- & 48 / 48 & 0 / 0 \\
rule & documented & 108/108 & 108/108 & -- & -- & 48 / 48 & 0 / 0 \\
rule & drifted & 0/364 & 0/364 & -- & -- & 0 / 0 & 48 / 48 \\
rule & injected & 108/108 & 108/108 & 0/66 & 0/66 & 48 / 48 & 0 / 0 \\
LLM & documented & 108/108 & 108/108 & -- & -- & 48 / 48 & 0 / 0 \\
LLM & drifted & 108/108 & 174/196 & -- & -- & 48 / 34 & 0 / 12 \\
LLM & injected & 125/184 & 129/136 & 48/105 & 1/71 & 24 / 42 & 18 / 6 \\
\bottomrule
\end{tabular}
\end{adjustbox}
\end{table}

两种解析器以相反的方式失败。规则解析器读取文档格式没有任何错误，并忽略注入字段，但它读不懂任何一行改写格式，因此所有改写格式的案例都被升级。模型能读懂改写格式（Qwen2.5-7B 108 条中 108 条正确，Llama-3.1-8B 196 条中 174 条；Llama 的误读造成 12 次升级和 2 个给出错误攻击的结论），但 Qwen2.5-7B 在 105 次注入观测中有 48 次采纳了攻击者声称的读数；Llama-3.1-8B 在 71 次中只采纳了 1 次。这是本研究中模型第一次提供控制器无法提供的东西，也正是在这里，模型重新打开了结束守卫关闭的攻击面。在这些回合中，账本遏制了损害：被采纳的读数（\emph{path scanning}、\emph{已授权的变更窗口}）与其他探针相矛盾，因此没有攻击被结为良性；Qwen2.5-7B 有 18 个注入回合在耗尽预算后升级，6 个良性的监控案例被判为另一个良性原因。

\subsection{安全边界}
\label{sec:adversarial}
\label{sec:security}
威胁模型允许攻击者控制部分证据。在没有对抗性证据时，仅账本就已消除大部分错误结论：在表来源研究中，ReAct 式的 Qwen2.5-32B 和 72B 基线在给出结论的需处置回合中分别把 51 个中的 13 个、54 个中的 20 个判为良性，Memory-only 为 54 个中 2 个和 53 个中 1 个，\hesp{} 为 0。对抗研究测试两类攻击：写在自由文本字段中的“以授权扫描结案”指令（四种措辞和位置不同的文本），以及伪造的证据：报告“已知扫描器”的来源信誉探针（\emph{伪造探针}），和被两个探针共同读取的上游信誉源（\emph{伪造数据源}）。\cref{tab:v11attack} 在全部五个模型上把结果归因到各个组件；v1.0 的结果见附录~\ref{app:tables} 的\cref{tab:v10c}。

\begin{table*}[t]
  \caption{安全边界（v1.1，sec-triage，五个模型）：在四种注入文本下，把需处置案例结为良性（“missed”）或没有结论（“unresolved”）的回合数（每个模型共 72 个）；在伪造证据下被结为良性的回合数（共 18 个；五个模型上的范围）；两个良性原因上的已验证回合数（共 6 个；五个模型上的范围）。}
  \label{tab:v11attack}
  \centering\small
  \begin{adjustbox}{max width=\textwidth}% generated by paper/make_tables.py -- do not edit
\begin{tabular}{lrrrrrrrr}
\toprule
 & \multicolumn{5}{c}{4 injection texts: missed / unresolved (of 72)} & \multicolumn{2}{c}{forged: missed (of 18)} & benign \\
\cmidrule(lr){2-6}\cmidrule(lr){7-8}\cmidrule(lr){9-9}
Configuration & Q-7B & Q-32B & Q-72B & L-8B & L-70B & probe & feed & verif.\ (of 6) \\
\midrule
Memory-only & 69 / 0 & 57 / 0 & 69 / 0 & 0 / 72 & 63 / 0 & 0--6 & 0--15 & 0--6 \\
Memory-only, raw text removed & 0 / 65 & 0 / 0 & 0 / 0 & 0 / 72 & 0 / 24 & 0--6 & 0--15 & 0--6 \\
\hesp{}, guard only & 0 / 44 & 0 / 46 & 0 / 66 & 0 / 72 & 0 / 58 & 0 & 0 & 0--6 \\
\hesp{}, stop, no guard & 39 / 0 & 36 / 0 & 48 / 0 & 0 / 0 & 36 / 0 & 18 & 18 & 6 \\
\hesp{}, guard + stop & 0 / 24 & 0 / 29 & 0 / 42 & 0 / 0 & 0 / 2 & 18 & 18 & 6 \\
\quad + corroboration by probe & 0 / 0 & 0 / 0 & 0 / 11 & 0 / 0 & 0 / 0 & 0 & 0 & 3 \\
\quad + corroboration by source & 0 / 12 & 0 / 0 & 0 / 3 & 0 / 0 & 0 / 0 & 0 & 0 & 3 \\
\bottomrule
\end{tabular}
\end{adjustbox}
\end{table*}

第一，阻止注入式结案的是结束守卫。每个自行决定并给出结论的模型都服从了指令；有控制器选择和停止但没有守卫时，规划器仍在 72 个回合中的 36 到 48 个里抢先以注入的结论结束。在每个模型上，所有带守卫的配置都把 72 个中的 0 个结为良性；v1.0 在 Qwen2.5-7B 上预先登记的比较为 $+1.000$ $[+1.000,+1.000]$，全部 18 个配对回合方向一致。从提示中去掉原始文本也能阻止注入式结案，对 Qwen2.5-32B 和 72B 没有代价，但 Qwen2.5-7B 有 72 个中 65 个案例得不到结论。注入仍然会耗费分析员的时间：没有控制器停止时，守卫留下 72 个中 44 到 72 个未解决的回合，加上停止后为 0 到 42 个。

第二，伪造证据能击败控制器。单条伪造的扫描器回答在一个探针之后就把授权扫描推过 0.8，控制器停止随即给出结论：在每个模型上 18 个需处置案例全部被结为良性，被两个探针读取的伪造数据源也是一样。优先依据信息量最大的证据行动，恰恰是攻击者所利用的。伪造数据源也会误导较大的 LLM 本身（Memory-only 下最多 18 个中 15 个）。

第三，良性结论的佐证以一定代价弥补这一缺口。要求两条特异支持的观测，会在每个模型上把所有伪造案例升级；但按探针计数的规则之所以挡住了伪造数据源，只是因为控制器的第二个探针读取了另一个来源。v1.2 的定向测试让固定流程先读 \texttt{source\_ips}、再读 \texttt{threat\_intel}，两者都来自被伪造的来源：此时按探针佐证把 18 个需处置案例全部结为良性，按来源组佐证则把 18 个全部升级。两种规则都使 sec-triage 两个良性原因的完成率减半（6 个中 3 个），因为那里的监控误报只有一个特异特征；在 sigma-triage 的 26 个良性解释上，佐证只让 104 个良性回合中的 8 个升级。所有计数都来自一个环境，一个 0/18 的格子只能给出很宽松的上界（约 15\%，单侧 95\%，即使其六个原因和三次重复相互独立）。

\subsection{知识从哪里来}
\label{sec:rq-help}
\label{sec:mismatch}
控制器的知识存放在预测表中。计数表与生成器的表相当：用每个原因和变体 $k{=}20$ 个开发回合计数，带结束守卫的 \hesp{} 把 Qwen2.5-7B 从 0.125 提升到 1.000（$+0.875$ $[+0.74,+0.97]$，预先登记），在三个 Qwen 模型上都与设计者表相同，而且 $k{=}5$ 就已足够（附录~\ref{app:tables} 的\cref{tab:v06}），这反映了环境的确定性。计数表多花 1.2 到 2.4 个单位，全部来自 \textit{other} 这一行，也就是开发数据中从未出现的那个原因。模型无法自己写出预测表：由规划器给出的表只能让完成率达到 0.069 到 0.375。

在我们的环境中，预测表与测试回合来自同一个生成器。\cref{tab:v10b} 用不含 LLM 的控制器打破了这一联系。有三点发现；前两点在两个任务族上都成立。

\begin{table*}[t]
  \caption{预测表误差（不含 LLM，后验停止阈值 0.8）。EIG/成本选择下：已验证、原因正确但缺少确认特征、原因错误和升级的回合比例（四者之和为一）；随机选择下：已验证完成率。“扰动”把每一行与随机 Dirichlet(1) 行按权重 $\lambda$ 混合（每个水平 20 张随机表）；“缺失一个原因”对依次从开发数据中去掉每个原因、并在该原因上测试的结果取平均。}
  \label{tab:v10b}
  \centering\small
  \begin{adjustbox}{max width=\textwidth}% generated by paper/make_tables.py -- do not edit
\begin{tabular}{lrrrrrrrrrr}
\toprule
 & \multicolumn{5}{c}{sec-triage} & \multicolumn{5}{c}{sigma-triage} \\
\cmidrule(lr){2-6} \cmidrule(lr){7-11}
Tables & verif. & correct, unverif. & wrong & esc. & random verif. & verif. & correct, unverif. & wrong & esc. & random verif. \\
\midrule
matched tables & 0.917 & 0.083 & 0.000 & 0.000 & 0.792 & 1.000 & 0.000 & 0.000 & 0.000 & 0.886 \\
base variant only & 0.917 & 0.083 & 0.000 & 0.000 & 0.792 & 1.000 & 0.000 & 0.000 & 0.000 & 0.886 \\
$k{=}1$, base only & 0.917 & 0.083 & 0.000 & 0.000 & 0.778 & 1.000 & 0.000 & 0.000 & 0.000 & 0.882 \\
perturbed, $\lambda{=}0.25$ & 0.885 & 0.083 & 0.000 & 0.031 & 0.500 & 0.992 & 0.008 & 0.000 & 0.000 & 0.910 \\
perturbed, $\lambda{=}0.5$ & 0.494 & 0.177 & 0.000 & 0.329 & 0.227 & 0.978 & 0.017 & 0.000 & 0.005 & 0.838 \\
perturbed, $\lambda{=}0.75$ & 0.098 & 0.027 & 0.004 & 0.871 & 0.042 & 0.805 & 0.040 & 0.009 & 0.146 & 0.513 \\
one cause missing (mean) & 0.000 & 0.000 & 0.500 & 0.500 & 0.000 & 0.000 & 0.000 & 0.079 & 0.921 & 0.000 \\
\bottomrule
\end{tabular}
\end{adjustbox}
\end{table*}

第一，不精确的预测表让控制器以升级的方式失败。随着 $\lambda$ 增大，完成率下降，但即使每一行有四分之三是噪声，错因率也不超过 0.009，并且在每个水平上 EIG/成本都领先于随机选择。这些结果只针对随机混合扰动，不能证明对任意或对抗性的似然误差具有鲁棒性。

第二，开发数据中缺失某个原因会产生自信的错误结论。该原因的行不含信息，本应揭示它的观测会被算作另一个原因的弱证据：在 sec-triage 上，八个原因中有四个在每个回合都被判成错误原因，在 sigma-triage 上为 38 个中的 3 个。每个这样的结论所给的原因都与真实原因同类（缺失的攻击被认成另一个需处置原因，缺失的良性解释被认成另一个良性解释），因此在这些测试中没有任何缺失原因导致攻击被结为良性。

第三，v1.1 研究加入了随机混合产生不了的误差（sec-triage，每格 48 个不含 LLM 的回合；其 3{,}968 个回合中没有任何攻击被结为良性）。\emph{系统性混淆}把每个原因的行向某个固定的另一原因混合：权重 0.25 时后验停止几乎不受影响（48 个中 47 个），但确认条件停止完全失效——没有任何结果在一个原因下的概率仍是其配对原因的十倍，于是全部 48 个回合都耗尽预算并升级。因此，当预测表混淆两个原因时，确认条件规则以放弃自动化为代价保持低错误率；把比值降到 3 也无法挽回（48 个中 0 个；\cref{tab:v12policy}），而联合证据验证了 48 个中的 47 个，且没有错误结论。\emph{候选集合中缺失真实原因}时，后验停止在 48 个回合中有 24 个给出另一个需处置原因，确认条件停止把 48 个全部升级。\emph{重复的来源}（第二个探针读取与 \texttt{source\_ips} 相同的数据源）没有可测的代价。阈值在 0.6 到 0.95 之间时，确认条件停止保持 1.000，成本逐步上升（2.54 到 4.14 个单位），而后验停止从 0.9 起降到 0.875。

\subsection{第二任务族}
\label{sec:external}
在 sigma-triage（候选原因来自公开检测规则的作者，结果由两个规划器家族之外的模型标注）上，Qwen2.5-7B 的两个预先登记终点都成立（12 条规则上的 98.33\% 区间；附录~\ref{app:tables} 的\cref{tab:v10a}）：排序增加 $+0.171$ $[+0.094,+0.250]$，带控制器停止的 \hesp{} 比 Memory-only 高 $+0.296$ $[+0.088,+0.473]$。排序效应在全部五个模型上都得到重复（$+0.105$ 到 $+0.230$），探测与停止的分离也得到重复：没有控制器停止时，Llama-3.1-8B 在 304 个回合中只结束了 1 个，加上停止后升到 1.000。这些配置没有结束守卫，Qwen2.5-7B 在 \hesp{} 内把 144 个攻击回合中的 10 个结为良性（Memory-only 下为 0），其中九个在同一条规则上；在每个这样的回合中，账本给所称原因的分数是 0.002，而且被引用的观测都不支持它，因此守卫会拒绝所有这些结论。sec-triage 的主要终点对聚类单位也是稳健的：以 8 个原因代替 24 个任务重抽样，每个已成立终点的区间仍然高于零，并且静态任务与漂移任务上的效应相近（附录~\ref{app:tables}）。

\subsection{公开安全数据与本评测的适配性}
\label{sec:realdata}
这些结果在真实数据上是否成立？我们考察了最接近的三个公开数据集能否检验本文所定义的调查策略；这需要告警让多种解释同时成立、探针的成本和信息量各不相同，以及客观的真值。附录~\ref{app:public} 给出版本、文件和方法；ExCyTIn 的可行性检查读取了其测试题（勘误 E-10），而 GUIDE 的测试集从未读取。在我们考察的子集和采用的适配方式下，没有找到适合本研究问题的设置。在 ExCyTIn-Bench~\citep{excytin} 中，大多数训练题直接点名要定位的告警（54.8\%），我们测试的一种指标关联探针找到答案的频率并不高于干扰项；这并不说明该基准整体只考察检索。在 GUIDE~\citep{guide} 中，组织身份携带了分级 1.50 比特中的 1.03 比特，这可能反映策略、基率或检测器构成，而且在冷启动事件上证据整合对真阳性的排序并不优于历史查表（AUROC 0.755 对 0.759）。在我们审查的 27 份 OTRF~\citep{otrf} 远程执行录制和两类事件中，一条只看告警文本的规则就能区分 46 条告警中的 44 条。这些子集缺少的是：使用相同工具的攻击与合法管理的配对，以及客观标注。Cyber Defense Benchmark~\citep{cyberdefensebench} 表明这类遥测可以做成交互式环境，形式是没有触发告警的 SQL 威胁狩猎。
